\title[Mass--Spring Oscillator]{Mass--Spring Oscillator}
\subtitle{Determining the spring constant dynamically · \labversionlabel}
\author{Mechanics Course}
\institute{Demonstration Material}
\date{2026}

\begin{document}

\begin{frame}[plain]
  \titlepage
\end{frame}

\begin{frame}{Experimental Question}
  \slidetag{OBJECTIVE}
  \vspace{0.7em}
  \begin{columns}[T,onlytextwidth]
    \begin{column}{0.57\textwidth}
      {\Large How does the period $T$ of vertical oscillations depend on the suspended mass $m$?}
      \vspace{0.9em}
      \begin{block}{Testable hypothesis}
        The squared period $T^2$ depends linearly on the mass.
      \end{block}
    \end{column}
    \begin{column}{0.38\textwidth}
      \resultcard{Result}{Spring constant $k$}
    \end{column}
  \end{columns}
\end{frame}

\begin{frame}{Equipment and Safety}
  \begin{columns}[T,onlytextwidth]
    \begin{column}{0.43\textwidth}
      \slidetag{EQUIPMENT}
      \vspace{0.55em}
      \begin{itemize}
        \item retort stand and spring;
        \item set of slotted masses;
        \item stopwatch;
        \item ruler.
      \end{itemize}
    \end{column}
    \begin{column}{0.52\textwidth}
      \begin{alertblock}{Before you begin}
        Check that the stand is secure. Support each mass while changing it,
        and do not stretch the spring beyond its elastic limit.
      \end{alertblock}
    \end{column}
  \end{columns}
\end{frame}

\begin{frame}{Physical Model}
  \keyformula{$m\ddot{x}+kx=0 \qquad T=2\pi\sqrt{\dfrac{m}{k}}$}
  \vspace{0.8em}
  \begin{itemize}
    \item oscillations have small amplitude;
    \item air resistance is negligible;
    \item the spring remains within its elastic range.
  \end{itemize}
\end{frame}

\begin{frame}{Apparatus Diagram}
  \begin{columns}[c,onlytextwidth]
    \begin{column}{0.42\textwidth}
      \centering
      \resizebox{0.72\linewidth}{!}{\begin{tikzpicture}[x=1cm,y=1cm]
  \fill[PhysicsPaper] (-2.0,-0.5) rectangle (2.4,5.0);
  \draw[physics line,line width=2pt] (-1.6,4.65) -- (1.6,4.65);
  \foreach \x in {-1.5,-1.2,...,1.5}
    \draw[PhysicsMuted,thin] (\x,4.65) -- ++(-0.22,0.22);
  \draw[physics accent,decorate,decoration={coil,aspect=0.45,segment length=5mm,amplitude=4mm}]
    (0,4.65) -- (0,1.75);
  \filldraw[fill=PhysicsBlue!18,draw=PhysicsNavy,line width=1.1pt,rounded corners=1mm]
    (-0.78,0.75) rectangle (0.78,1.75);
  \node[font=\bfseries\sffamily,text=PhysicsNavy] at (0,1.25) {$m$};

  \draw[physics force] (0.95,1.25) -- ++(0,1.25)
    node[right,physics label] {$\vec F_{\text{spring}}$};
  \draw[physics force] (-0.95,1.25) -- ++(0,-1.15)
    node[left,physics label] {$m\vec g$};
  \draw[-{Latex[length=3mm]},PhysicsBlue,line width=1.2pt]
    (1.8,2.2) -- (1.8,0.45) node[below,physics label] {$x$};
  \draw[PhysicsMuted,dashed] (-1.35,1.75) -- (1.55,1.75)
    node[right,physics label] {equilibrium};
\end{tikzpicture}
}
    \end{column}
    \begin{column}{0.53\textwidth}
      \begin{block}{Coordinate system}
        The positive $x$-axis points downward from the equilibrium position.
      \end{block}
      \vspace{0.5em}
      At equilibrium, the weight and spring force balance. After a
      displacement, the net force restores the mass toward equilibrium.
    \end{column}
  \end{columns}
\end{frame}

\begin{frame}{Measurement Procedure}
  \begin{enumerate}
    \item Attach a mass and wait until it reaches equilibrium.
    \item Displace the mass downward by no more than \SI{2}{\centi\metre}.
    \item Measure the time $t_N$ for $N=20$ complete oscillations.
    \item Repeat the measurement for five different masses.
    \item For each trial, calculate $T=t_N/N$ and $T^2$.
  \end{enumerate}
\end{frame}

\begin{frame}{Raw-Data Table}
  \centering
  \renewcommand{\arraystretch}{1.22}
  \begin{tabular}{@{}ccccc@{}}
    \toprule
    $m$, \si{\kilo\gram} & $N$ & $t_N$, \si{\second} & $T$, \si{\second} & $T^2$, \si{\second\squared} \\
    \midrule
    \phantom{0.000} & 20 & \phantom{00.00} & \phantom{0.000} & \phantom{0.000} \\
    & 20 & & & \\
    & 20 & & & \\
    & 20 & & & \\
    & 20 & & & \\
    \bottomrule
  \end{tabular}

  \vspace{0.9em}
  {\color{PhysicsMuted}\small Record raw measurements before performing any intermediate rounding.}
\end{frame}

\begin{frame}{Data Analysis}
  \keyformula{$T^2=am+b \qquad k=\dfrac{4\pi^2}{a}$}
  \vspace{0.8em}
  \begin{enumerate}
    \item Plot $T^2$ against $m$.
    \item Determine the slope $a$ of the linear fit.
    \item Calculate $k$ and discuss the physical meaning of the intercept $b$.
  \end{enumerate}
\end{frame}

\begin{frame}{Check-for-Understanding Tasks}
  \slidetag{SELF-CHECK}
  \vspace{0.6em}
  \begin{block}{1 · Units}
    Determine the SI unit of the slope $a$ on a graph of $T^2$ against $m$.
  \end{block}
  \begin{block}{2 · Calculation}
    An experiment gives $a=\SI{0.80}{\second\squared\per\kilo\gram}$.
    Determine the spring constant.
  \end{block}
\end{frame}

\begin{solutionframe}[Solution · Units of the Slope]
  From $T^2=am+b$,
  \[
    a=\frac{T^2}{m},
    \qquad
    [a]=\si{\second\squared\per\kilo\gram}.
  \]
  \answer{$\si{\second\squared\per\kilo\gram}$.}
\end{solutionframe}

\begin{solutionframe}[Solution · Spring Constant]
  \[
    k=\frac{4\pi^2}{a}
     =\frac{4\pi^2}{\SI{0.80}{\second\squared\per\kilo\gram}}
     \approx\SI{49}{\newton\per\metre}.
  \]
  \answer{$k\approx\SI{49}{\newton\per\metre}$.}
  \vspace{0.5em}
  \teacheronly{This illustrative value is not the result of an actual experiment.}
\end{solutionframe}

\begin{frame}{Report Requirements}
  \begin{columns}[T,onlytextwidth]
    \begin{column}{0.48\textwidth}
      \begin{itemize}
        \item objective and apparatus diagram;
        \item raw-data table;
        \item graph of $T^2$ against $m$;
      \end{itemize}
    \end{column}
    \begin{column}{0.48\textwidth}
      \begin{itemize}
        \item fitted coefficients $a$ and $b$;
        \item value of $k$ with its unit;
        \item conclusion addressing the hypothesis.
      \end{itemize}
    \end{column}
  \end{columns}
\end{frame}

\end{document}
